\appendix

\chapter{Anhang}
\label{ch:anhang}

\section{Reproduzierbarkeit der Experimente}
\label{sec:appendix-reproducibility}

Die in Abschnitt~\ref{sub-sec:lps-experimente} und Abschnitt~\ref{sub-sec:alps-experimentelle-evaluation} beschriebenen Experimente sind reproduzierbar, indem die benötigten Docker-Images geladen, eigene Konfigurationsdateien erstellt und die Simulationen anschließend gestartet werden.

\subsection{Voraussetzungen}

\begin{itemize}
    \item Installiertes Docker~\citep{docker_2023}.
    \item Installiertes Docker Compose~\citep{docker_compose}.
\end{itemize}

\subsection{Docker-Images laden}

Die folgenden Befehle laden alle benötigten Docker-Images.

\begin{bashcode}[title={Docker-Images laden}]
# RabbitMQ-Broker
docker pull rabbitmq:4-management
# Simulation-Runner und Kernkomponenten von ALPS
# (ALP I, ALP II, Alarmanlagen-Simulator, Server-Simulator, Validator, TCP-Proxy)
docker pull jadkhaddad/lps-rs
# ALP Dedup
docker pull jadkhaddad/lps-gleam
# Auswertungsskripte zur Plot-Erzeugung
docker pull jadkhaddad/lps-scripts
# Kernkomponenten von LPS
# (SRV I PROXY, SRV II PROXY, LP I, LP II)
docker pull jadkhaddad/lps-mqtt-python
# HiveMQ-Broker mit den verwendeten Plugins
docker pull jadkhaddad/lps-mqtt-hivemq
\end{bashcode}

\subsection{Docker-Netzwerk erstellen}

Vor dem Start wird ein Docker-Netzwerk angelegt, über das die Komponenten miteinander kommunizieren.

\begin{bashcode}[title={Docker-Netzwerk erstellen}]
docker network create lps-simulation-net
\end{bashcode}

\subsection{Konfigurationsdateien anlegen}

Vor dem Start werden die benötigten Simulationskonfigurationen als \gls{yaml}-Dateien angelegt. Diese Dateien enthalten die Parameter für Last, Fehlerprofile und Laufzeit der Experimente. Eine Konfigurationsdatei besteht aus einem allgemeinen Teil und einem szenariospezifischen Teil, entweder für \gls{alps} oder für \gls{lps}. Parameter, die mit \texttt{(optional)} gekennzeichnet sind, können weggelassen werden. Fehlerraten werden als Wahrscheinlichkeit zwischen \texttt{0} und \texttt{1} angegeben. Zeitwerte werden als Dauerangaben angegeben, z.\,B.~\texttt{500ms} oder \texttt{10s}.

\begin{yamlcode}[title={Allgemeiner Teil}]
# Zielverzeichnis der erzeugten Simulationsartefakte
simulation_dir: "simulation"
# Externes Docker-Netzwerk, in dem die Services gestartet werden
simulation_docker_network: "lps-simulation-net"

# Image für Simulation-Runner und Kernkomponenten von ALPS
lps_rust_image: "jadkhaddad/lps-rs"
# Image für die Auswertungsskripte
lps_scripts_image: "jadkhaddad/lps-scripts"
# Image für ALP Dedup
lps_gleam_image: "jadkhaddad/lps-gleam"
# Image für die die Kernkomponenten von LPS
lps_mqtt_python_image: "jadkhaddad/lps-mqtt-python"
# Image für den HiveMQ-Broker
lps_mqtt_hivemq_image: "jadkhaddad/lps-mqtt-hivemq"

# Simulierte Alarmanlagen
devices:
  # Anzahl simulierter Alarmanlagen
  devices_count: 100
  # Maximales Intervall zwischen zwei Meldungen pro Alarmanlage
  max_event_interval: 500ms
  # (optional) Minimales Intervall zwischen zwei Meldungen
  min_event_interval: 100ms
  # Maximale Anzahl erzeugter Meldungen pro Alarmanlage
  max_events: 100
  # Wartezeit bis zum erneuten Sendeversuch nach einem Fehler
  retry_duration: 500ms
# Server-Simulator
server:
  # Adresse, unter der der Server-Simulator erreichbar ist
  socket_addr: "0.0.0.0:3000"
  # (optional) Nach außen freigegebene Adresse des Servers
  expose_socket_addr: "0.0.0.0:3000"
  # Anzahl gestarteter Server-Instanzen (mindestens 1)
  server_count: 3
  # (optional) Aktiviert ALP Dedup
  deduplication: true
  # (optional) Server beendet sich nach dieser Leerlaufzeit, wenn keine Anfragen mehr eintreffen
  shutdown_after_idle_duration: 30s

  # Zufallsbasiertes Verhalten (ignoriert bei use_gs_prod_workload: true)

  # Wahrscheinlichkeit für Fehlerantworten
  failure_rate: 0.1
  # (optional) Maximale Bearbeitungszeit einer Anfrage
  max_work_duration: 50ms
  # (optional) Minimale Bearbeitungszeit einer Anfrage
  min_work_duration: 10ms
  # (optional) Maximale zusätzliche Antwortverzögerung
  max_response_duration: 50ms
  # (optional) Minimale zusätzliche Antwortverzögerung
  min_response_duration: 10ms

  # (optional) Bearbeitungszeiten werden vom Vorhersagemodell bestimmt
  use_gs_prod_workload: false
  # (optional) Hintergrundlast in Anfragen pro Sekunde
  # (nur wirksam bei use_gs_prod_workload: true)
  background_load_per_second: 10.0
\end{yamlcode}

\newpage

\begin{yamlcode}[title={ALPS-spezifischer Teil}]
amqp:
  # Anzahl der Broker-Queues
  queues_count: 10
  producers:
    # Anzahl der ALP I-Instanzen (mindestens 1)
    producers_count: 3
    # Adresse der HTTP-Schnittstelle von ALP I
    socket_addr: "0.0.0.0:8080"
  consumers:
    # Anzahl der ALP II-Instanzen (mindestens 1)
    consumers_count: 3
    # Wartezeit vor dem Reconnect nach einem Verbindungsabbruch
    reconnect_duration: 500ms
  # (optional) Fehlerproxy zwischen ALP II und Broker
  amqp_proxy:
    # Adresse des Fehlerproxys
    socket_addr: "0.0.0.0:5673"
    # Wahrscheinlichkeit eines künstlichen Verbindungsfehlers
    failure_rate: 0.1
    # Maximale Dauer einer Verbindung vor erzwungener Trennung
    max_up_duration: 10s
    # (optional) Minimale Dauer einer Verbindung
    min_up_duration: 1s
  # (optional) Fehlerproxy zwischen ALP II und Server
  server_proxy:
    socket_addr: "0.0.0.0:3001"
    failure_rate: 0.1
    max_up_duration: 10s
    min_up_duration: 1s
\end{yamlcode}

\newpage

\begin{yamlcode}[title={LPS-spezifischer Teil}]
mqtt:
  # (optional) Verwendet Mosquitto statt HiveMQ als Broker
  use_mosquitto: false
  producers:
    # Anzahl der LP I-Instanzen (mindestens 1)
    producers_count: 3
  consumers:
    # Anzahl der LP II-Instanzen (mindestens 1)
    consumers_count: 3
  # (optional) Fehlerproxy zwischen LP II und Broker
  mqtt_proxy:
    socket_addr: "0.0.0.0:1884"
    failure_rate: 0.1
    max_up_duration: 10s
    min_up_duration: 1s
\end{yamlcode}

\subsection{Simulationen starten}

Die Simulationen werden aus dem Verzeichnis gestartet, in dem die Konfigurationsdateien liegen.

\begin{bashcode}[title={ALPS-Simulation starten}]
docker run --rm -it \
-v /var/run/docker.sock:/var/run/docker.sock \
-v "$(pwd):/app" \
    jadkhaddad/lps-rs \
    simulation-runner \
    --config-file <amqp-config.yaml> \
    --root-abs-path $(pwd)
\end{bashcode}

\begin{bashcode}[title={LPS-Simulation starten}]
docker run --rm -it \
-v /var/run/docker.sock:/var/run/docker.sock \
-v "$(pwd):/app" \
    jadkhaddad/lps-rs \
    mqtt-simulation-runner \
    --config-file <mqtt-config.yaml> \
    --root-abs-path $(pwd)
\end{bashcode}

\subsection{Ergebnisse}

Nach Abschluss eines Laufs befinden sich die erzeugten Ergebnisse im Verzeichnis  \texttt{\$(pwd)/<simulation\_dir>/mount}

\section{Quellcode}
\label{sec:source-code}

Dieses Projekt ist ein Open-Source-Projekt~\citep{source} und unter der \textit{MIT License}~\citep{MIT} lizenziert.

\begin{figure}[h]
    \centering
    \includegraphics[width=0.4\textwidth]{./qrcode/qrcode.pdf}
    \caption{Quellcode}
    \label{fig:source-code}
\end{figure}